\section{Executive summary}
\label{sec:executive-summary}

\begin{notebox}
\noteboxtitle{Question answered by this report}
If PR \#1811 changes how Lido accounts for deposits, reserves, module balances,
and rewards, can the selected accounting rules be stated precisely enough that a
machine can check them in a reproducible focused model? This report package
gives an executable local answer for the six selected P0 (highest-priority)
properties over the focused economic model, under the named assumptions in
\cref{sec:trust-boundary}.
\end{notebox}

\noindent In this repository, ``executable'' has a specific meaning: the model
target is listed, all premises not discharged inside the model are registered by
assumption ID, and \texttt{make test} plus \texttt{make prove} recheck the target
set from a clean checkout.

\subsection*{The story in one page}

The core safety question is simple enough to ask without Solidity or Lean:
\emph{when ETH and validator balances move through SRv3 accounting, can any
value be spent from the wrong place, counted twice, rewarded from stale data, or
sent through a disabled path?} The report answers that question by turning each
informal concern into a small conservation law.

\smallskip
\noindent A useful mental model is three ledgers that must agree:

\begin{enumerate}
  \item The \textbf{buffer ledger} says which ETH is depositable and which ETH is
  reserved for withdrawals.
  \item The \textbf{router ledger} says how much ETH is pulled for deposits and
  how validator balances are summed across staking modules.
  \item The \textbf{reward ledger} says which accepted balances may be used to
  compute rewards, fees, and recipients.
\end{enumerate}

\noindent An ``accepted report'' means oracle-provided module-balance data that
has passed the modeled SRv3 validation checks and is therefore allowed to update
accounting state. The proof lane checks that PR \#1811 preserves the three
ledgers through those modeled transitions. It is not trying to prove every
operational fact about deployment, cryptography, governance, gas, or the
consensus layer.

\begin{notebox}
\noteboxtitle{Why this is more than testing}
A test asks, ``does this example work?'' A Verity/Lean proof target asks,
``does this rule follow from the modeled transition for all values covered by
the theorem statement?'' That is why the report spends so much effort naming the
model and assumptions: they define exactly what the local command checks.
\end{notebox}

\smallskip
\noindent The executable claim this report package supports is:

\begin{quote}
For the focused SRv3 economic state machine, the Verity/Lean model preserves
deposit conservation, reserve separation, module-balance conservation,
report-before-reward consistency, reward/fee bounds, and flow-specific status
gating under the named assumptions.
\end{quote}

\noindent This is not a claim that every deployed-system behavior is proved. The
exact boundary is part of the result: cryptographic witnesses, consensus
truthfulness, governance configuration, and deployment refinement are named in
\cref{sec:trust-boundary}.

\subsection*{How to read the proof}

Read each row below as a three-step trace. First, the
Verity/Lean model names the transition being studied. Second, the property
states what must not go wrong in the focused model. Third, Lake checks the
theorem file against the explicit assumptions.

\smallskip
\noindent The matrix answers six plain questions: what ETH is spendable, what
deposits consume, what balances add up to, what report rewards use, how fees are
bounded, and which module states block economic flow.

\subsection*{Proof matrix}

\noindent All six rows below have executable artifacts in \path{proofs/},
\path{LidoSRv3/}, and \path{verity/targets/}.

\noindent
\renewcommand{\arraystretch}{1.18}
\begin{tabularx}{\linewidth}{@{}>{\raggedright\arraybackslash}p{0.12\linewidth}>{\raggedright\arraybackslash}p{0.24\linewidth}>{\raggedright\arraybackslash}X@{}}
  \toprule
  \thead{ID} & \thead{PR \#1811 surface} & \thead{Executable property} \\
  \midrule
  SRV3-P1 & Lido buffer allocation &
  Withdrawal-reserved ETH is excluded from the depositable resource, so deposit
  spending cannot borrow from withdrawal liquidity. \\
  \rowsep
  SRV3-P2 & Router deposit path &
  Router-pulled ETH equals 32 ETH times actual deposits and is fully sunk
  into CL deposits, with no modeled residual held by the router. \\
  \rowsep
  SRV3-P3 & Module balances &
  Router-wide validator balance is exactly the sum of accepted per-module
  balances, so the aggregate cannot drift away from the module ledger. \\
  \rowsep
  SRV3-P4 & Oracle report ordering &
  Module fee-distribution weights come from the current accepted module-balance
  state, not from stale or partially applied reports. \\
  \rowsep
  SRV3-P5 & Rewards and fees &
  Module fee-distribution weights and recipient alignment use accepted module
  balances; total reward and fee magnitude remain accepted report or
  vault-context inputs under assumptions. \\
  \rowsep
  SRV3-P6 & Status gates &
  Non-active or deposit-paused modules receive no modeled deposit allocation;
  stopped modules receive no module-side reward or fee allocation. \\
  \bottomrule
\end{tabularx}
\renewcommand{\arraystretch}{1.0}

\subsection*{Review posture}

This report package is written for two readers at once. A protocol reviewer can follow
the ledgers and the trust boundary without reading code. A formal-methods
reviewer can check whether the target IDs, assumptions, and rebuild gates support
the advertised local executable claim.
